% ============================================================================
% MUSTERLÖSUNG DS 7 - Grupo A + B: Cultura y Diversión
% ============================================================================

\documentclass[11pt,a4paper]{article}

\usepackage[utf8]{inputenc}
\usepackage[T1]{fontenc}
\usepackage[ngerman]{babel}
\usepackage[left=2cm,right=2cm,top=1.8cm,bottom=1.5cm]{geometry}
\usepackage{helvet}
\usepackage{xcolor}
\usepackage{tabularx}
\usepackage{array}
\usepackage{enumitem}
\usepackage{tcolorbox}
\usepackage{fancyhdr}
\usepackage{lastpage}

\renewcommand{\familydefault}{\sfdefault}

\definecolor{spanisch}{HTML}{c41e3a}
\definecolor{grau}{HTML}{666666}
\definecolor{hellgrau}{HTML}{f5f5f5}
\definecolor{gruppeB}{HTML}{1565c0}
\definecolor{loesung}{HTML}{c62828}

\pagestyle{fancy}
\fancyhf{}
\fancyhead[L]{\small\textcolor{grau}{Spanisch Spätbeginner -- MUSTERLÖSUNG}}
\fancyhead[R]{\small\textcolor{grau}{ML-07}}
\fancyfoot[C]{\small\thepage/\pageref{LastPage}}
\renewcommand{\headrulewidth}{0.5pt}

\newcommand{\lsg}[1]{\textcolor{loesung}{\textbf{#1}}}
\newcommand{\es}[1]{\textit{#1}}

\newtcolorbox{gruppenbox}[1][]{
    colback=white,
    colframe=spanisch,
    fonttitle=\bfseries\color{white},
    colbacktitle=spanisch,
    title=#1,
    boxrule=1pt,
    arc=0pt,
    left=5pt, right=5pt, top=5pt, bottom=5pt
}

\newtcolorbox{gruppeBbox}[1][]{
    colback=white,
    colframe=gruppeB,
    fonttitle=\bfseries\color{white},
    colbacktitle=gruppeB,
    title=#1,
    boxrule=1pt,
    arc=0pt,
    left=5pt, right=5pt, top=5pt, bottom=5pt
}

\begin{document}

% ============================================================================
% KOPFBEREICH
% ============================================================================
\begin{center}
{\Large\textbf{\textcolor{loesung}{MUSTERLÖSUNG}}}\\[0.3em]
{\large DS 7 -- Cultura y Diversión}
\end{center}

\vspace{0.5em}
\hrule
\vspace{1em}

% ============================================================================
% GRUPPE A
% ============================================================================
\begin{gruppenbox}[Gruppe A: Cultura y Diversión (A2-Niveau)]

\textbf{Kasten 1: Países hispanohablantes (21 países)}

\textbf{Sudamérica (8 Länder):} \lsg{Argentina, Bolivia, Chile, Colombia, Ecuador, Paraguay, Perú, Uruguay, Venezuela}

(Hinweis: 9 Länder in Südamerika, davon 8 im Kasten einzutragen)

\vspace{0.8em}

\textbf{Kasten 2: Adjetivos para describir personas}

\begin{tabularx}{\textwidth}{@{}|X|X|X|X|@{}}
\hline
\textbf{positiv} & \textbf{Deutsch} & \textbf{negativ} & \textbf{Deutsch} \\
\hline
simpático/-a & \lsg{sympathisch} & antipático/-a & \lsg{unsympathisch} \\
\hline
inteligente & \lsg{intelligent} & tonto/-a & \lsg{dumm} \\
\hline
trabajador/-a & \lsg{fleißig} & perezoso/-a & \lsg{faul} \\
\hline
generoso/-a & \lsg{großzügig} & egoísta & \lsg{egoistisch} \\
\hline
paciente & \lsg{geduldig} & impaciente & \lsg{ungeduldig} \\
\hline
\end{tabularx}

\vspace{0.8em}

\textbf{Aufgabe 1: Países y capitales (6 P.)}

\begin{tabularx}{\textwidth}{@{}|X|X|@{}}
\hline
\textbf{País} & \textbf{Capital} \\
\hline
España & \lsg{Madrid} \\
\hline
México & \lsg{Ciudad de México} \\
\hline
Argentina & \lsg{Buenos Aires} \\
\hline
Colombia & \lsg{Bogotá} \\
\hline
Perú & \lsg{Lima} \\
\hline
Chile & \lsg{Santiago de Chile} \\
\hline
\end{tabularx}

\vspace{0.5em}

\textbf{Aufgabe 2: Adjektive ergänzen (6 P.)}
\begin{enumerate}[label=\alph*), leftmargin=2em, nosep]
    \item Mi hermano nunca estudia. Es muy \lsg{perezoso}.
    \item Elena siempre ayuda a sus amigos. Es muy \lsg{simpática / generosa}.
    \item Pablo trabaja mucho y nunca descansa. Es muy \lsg{trabajador}.
    \item María siempre sonríe y es agradable. Es muy \lsg{simpática}.
    \item Juan comparte todo con los demás. Es muy \lsg{generoso}.
    \item Ana nunca espera. Es muy \lsg{impaciente}.
\end{enumerate}

\vspace{0.5em}

\textbf{Aufgabe 3: Lied ``Vivir mi vida'' (8 P.)}

\begin{tabularx}{\textwidth}{@{}|X|@{}}
\hline
Voy a \lsg{reír}, voy a bailar \\
Vivir mi \lsg{vida}, la la la la \\
Voy a \lsg{bailar}, voy a gozar \\
Vivir mi \lsg{vida}, la la la la \\[0.3em]
A veces llega la \lsg{lluvia} \\
A veces me \lsg{siento triste} yo también \\
Mas yo sé bien / Que hay mañana \\
\hline
\end{tabularx}

\vspace{0.5em}

\textbf{Aufgabe 4: ¿Quién soy? (8 P.) -- Beispiellösung für Lionel Messi:}

\lsg{1. Es un hombre.}\\
\lsg{2. Es de Argentina.}\\
\lsg{3. Es deportista.}\\
\lsg{4. Juega al fútbol.}\\
\lsg{5. Es muy famoso en todo el mundo.}

\vspace{0.5em}

\textbf{Aufgabe 5: Mi país favorito (6 P.) -- Beispiellösung:}

\lsg{Mi país favorito es México. La capital es Ciudad de México. México está en Norteamérica. En México se habla español. México es famoso por su comida deliciosa, especialmente los tacos y las enchiladas.}

\end{gruppenbox}

\vspace{1em}

% ============================================================================
% GRUPPE B
% ============================================================================
\begin{gruppeBbox}[Gruppe B: Cultura y Diversión (B1-Niveau)]

\textbf{Kasten 1: El mundo hispanohablante -- Datos importantes}

\begin{tabularx}{\textwidth}{@{}|p{3cm}|X|@{}}
\hline
\textbf{Información} & \textbf{Datos} \\
\hline
Número de países & \lsg{21} países hispanohablantes \\
\hline
Hablantes totales & Aprox. \lsg{500--600} millones de personas \\
\hline
País más grande & \lsg{Argentina} (por superficie) \\
\hline
País más poblado & \lsg{México} (por habitantes) \\
\hline
Continentes & Europa, \lsg{América (Norte, Centro, Sur), Caribe}, África \\
\hline
\end{tabularx}

\vspace{0.8em}

\textbf{Aufgabe 1: Países y datos (8 P.)}

\begin{tabularx}{\textwidth}{@{}|p{2.5cm}|X|X|p{2.5cm}|@{}}
\hline
\textbf{País} & \textbf{Capital} & \textbf{Continente} & \textbf{Moneda} \\
\hline
España & \lsg{Madrid} & Europa & Euro \\
\hline
México & Ciudad de México & \lsg{Norteamérica} & \lsg{Peso mexicano} \\
\hline
Argentina & \lsg{Buenos Aires} & Sudamérica & Peso \\
\hline
Colombia & Bogotá & \lsg{Sudamérica} & \lsg{Peso colombiano} \\
\hline
Cuba & \lsg{La Habana} & Caribe & \lsg{Peso cubano} \\
\hline
Chile & Santiago & \lsg{Sudamérica} & Peso \\
\hline
Perú & \lsg{Lima} & Sudamérica & Sol \\
\hline
Venezuela & Caracas & \lsg{Sudamérica} & \lsg{Bolívar} \\
\hline
\end{tabularx}

\vspace{0.5em}

\textbf{Aufgabe 2: Adjektive mit Intensificadores (6 P.)}
\begin{enumerate}[label=\alph*), leftmargin=2em, nosep]
    \item \lsg{Shakira es muy talentosa.}
    \item \lsg{Rafael Nadal es bastante trabajador.}
    \item \lsg{Pablo Picasso es muy creativo.}
    \item \lsg{Frida Kahlo es bastante original.}
    \item \lsg{Penélope Cruz es muy guapa.}
    \item \lsg{Gabriel García Márquez es muy inteligente.}
\end{enumerate}

\vspace{0.5em}

\textbf{Aufgabe 3: Personenbeschreibung raten (6 P.)}
\begin{enumerate}[label=\alph*), leftmargin=2em, nosep]
    \item \lsg{Lionel Messi}
    \item \lsg{Shakira}
    \item \lsg{Rafael Nadal}
\end{enumerate}

\vspace{0.5em}

\textbf{Aufgabe 4: Beschreibung einer berühmten Person (8 P.) -- Beispiellösung:}

\lsg{Frida Kahlo es una artista mexicana muy famosa. Es bastante original y muy creativa. Sus pinturas son un poco tristes pero muy interesantes. Es famosa por sus autorretratos y por su estilo único. Frida Kahlo es una de las artistas más importantes de México.}

\vspace{0.5em}

\textbf{Aufgabe 5: Comparación cultural (6 P.) -- Beispiellösung:}

\lsg{He elegido México. Mientras que en Alemania el clima es bastante frío, en México hace mucho calor. A diferencia de Alemania, en México se habla español. También hay diferencias en la comida: en México comen mucho picante, en cambio en Alemania la comida es menos picante.}

\end{gruppeBbox}

\vspace{1em}
\hrule
\vspace{0.3em}

\textbf{Bewertungshinweise:}
\begin{itemize}[leftmargin=1.5em, nosep]
    \item Gruppe A: 34 Punkte gesamt
    \item Gruppe B: 34 Punkte gesamt
    \item Bei freien Texten: Kommunikation wichtiger als perfekte Grammatik
    \item Adjektive: Auch sinnvolle Alternativen akzeptieren
    \item Länder-Fakten: Aktuelle Informationen akzeptieren (Währungen können sich ändern)
    \item Halbe Punkte möglich bei teilweise richtigen Antworten
\end{itemize}

\end{document}
